\hypertarget{appendix-48d2e66378ab}{}
\section*{Matrix inversion pivoting and residuals}
\section{Matrix inversion pivoting and residuals}\label{matrix-inversion-pivoting-and-residuals}

\textbf{Curriculum depth:} Introduced · \textbf{First seen:} Phase7 · \textbf{Latest development:} Phase7

\subsection{Definition and mechanism}\label{definition-and-mechanism}

Inversion reconstructs input coordinates from output for an invertible square map. Pivoted elimination tracks row operations on an augmented system. A residual checks a computed answer; conditioning determines how strongly input errors can affect it. Prefer solving a system when an explicit inverse is unnecessary.

\textbf{Reconstruction cue:} Check both A times computed inverse and the inverse times A against identity, with scale-aware residuals.

\subsection{Prerequisites}\label{prerequisites}

\begin{itemize}
\tightlist
\item
  \hyperref[appendix-da478be01d02]{Determinants orientation and volume}
\item
  \hyperref[appendix-8f6a01f284a7]{Conditioning and singularity}
\end{itemize}

\subsection{Used by}\label{used-by}

\begin{itemize}
\tightlist
\item
  \href{}{Coordinate frames and change of basis (02-Foundations-II, appendix-1bae9457dbc3)}
\end{itemize}

\subsection{Recall and next depth}\label{recall-and-next-depth}

Reconstruct the example, state its conventions and find an input that violates its preconditions. Deeper algorithms and measured production performance remain separate from the present model.

\subsection{Source and navigation}\label{source-and-navigation}

\href{Continuous\%20Notes/Computing\%20Foundations\%20-\%20From\%20Bits\%20to\%20C++\%20and\%20Systems\%20APIs\#Keep\%20the\%20promises\%20visible\%20while\%20moving\%20between\%20representations}{Continuous explanation} · \hyperref[appendix-fd090d982a15]{Coverage and source gap} · \hyperref[appendix-a103e7846d21]{Sources and verification}
